\subsection{Allgemeine Korrespondenzen}
%hier können wir beispielhaft Korrespondenzen zu Zeitfunktionen im Fregu herleiten

\begin{table}[H]
    \centering
    \resizebox{\textwidth}{!}{%
    \begin{tabular}{|c|c||c|c|}
    \hline
    \textbf{Zeitfunktion} $f(t)$ & \textbf{Transformierte} $F(s)$ & \textbf{Zeitfunktion} $f(t)$ & \textbf{Transformierte} $F(s)$ \\
    \hline
    $1$ & $\frac{1}{s}$ & $\sinh^2(at)$ & $\frac{2a^2}{s(s^2 - 4a^2)},\ \text{Re}(s) > 2|a|$ \\
    $\exp(at)$ & $\frac{1}{s - a},\ \text{Re}(s) > a$ & $\cosh^2(at)$ & $\frac{s^2 - 2a^2}{s(s^2 - 4a^2)},\ \text{Re}(s) > 2|a|$ \\
    $\frac{t^n}{n!}\exp(at)$ & $\frac{1}{(s - a)^{n+1}},\ \text{Re}(s) > a$ & $t \sin(\omega t)$ & $\frac{2\omega s}{(s^2 + \omega^2)^2}$ \\
    $\cos(\omega t)$ & $\frac{s}{s^2 + \omega^2}$ & $t \cos(\omega t)$ & $\frac{s^2 - \omega^2}{(s^2 + \omega^2)^2}$ \\
    $\sin(\omega t)$ & $\frac{\omega}{s^2 + \omega^2}$ & $\exp(-\delta t) \cos(\omega t)$ & $\frac{s + \delta}{(s + \delta)^2 + \omega^2},\ \text{Re}(s) + \delta > 0$ \\
    $\sin^2(\omega t)$ & $\frac{s}{2\omega^2}\cdot \frac{1}{s^2 + 4\omega^2}$ & $\exp(-\delta t) \sin(\omega t)$ & $\frac{\omega}{(s + \delta)^2 + \omega^2},\ \text{Re}(s) + \delta > 0$ \\
    $\cos^2(\omega t)$ & $\frac{s^2 + 2\omega^2}{s(s^2 + 4\omega^2)}$ & $\frac{\delta \sin(\omega t) - \omega \sin(\delta t)}{\omega(\delta^2 - \omega^2)}$ & $\frac{1}{(s^2 + \omega^2)(s^2 + \delta^2)},\ \omega \ne \delta$ \\
    $\sin(\omega t + \varphi)$ & $\frac{s \sin(\varphi) + \omega \cos(\varphi)}{s^2 + \omega^2}$ & $\cos(\omega t) - \cos(\delta t)$ & $\frac{s}{(s^2 + \omega^2)(s^2 + \delta^2)},\ \omega \ne \delta$ \\
    $\cos(\omega t + \varphi)$ & $\frac{s \cos(\varphi) - \omega \sin(\varphi)}{s^2 + \omega^2}$ & $\sin(\omega t) - \omega t \cos(\omega t)$ & $\frac{1}{(s^2 + \omega^2)^2},\ \omega \ne 0$ \\
    $t^n$ & $\frac{n!}{s^{n+1}}$ & $\sin(\omega t) + \omega t \cos(\omega t)$ & $\frac{2\omega s}{(s^2 + \omega^2)^2},\ \omega \ne 0$ \\
    $\cosh(at)$ & $\frac{s}{s^2 - a^2},\ \text{Re}(s) > |a|$ & $\delta(t)$ (siehe Seite 898) & $1$ \\
    $\sinh(at)$ & $\frac{a}{s^2 - a^2},\ \text{Re}(s) > |a|$ & & \\
    \hline
    \end{tabular}}
    \caption{Tabelle der allgemeinen Laplace-Transformierten \cite[S.~883]{buch1}}
    \label{tab:laplace_allgemein}
    \end{table}
    
    
Beispielsweise wollen wir nun einmal den Term:

\[
f(t) = \sin^2(\omega t)
\]

Laplace transfomieren.
Zunächst stellen wir also unser Laplace Integral mit den Grenzen 0 und \(\infty\) auf:

\[
\mathcal{L} \{ \sin^2(\omega t) \} = \int_0^{\infty} e^{-st} \sin^2(\omega t) \, dt
\]

Zur Berechung der Stammfunktion nutzen wir folgendes Additionstheorem:

\[
\sin^2(\omega t) = \frac{1 - \cos(2\omega t)}{2}
\]

Eingesetzt in unser Integral können wir dieses nun in zwei zu addierende Integrale aufteilen:

\[
\int_0^{\infty} e^{-st} \cdot \frac{1 - \cos(2\omega t)}{2} \, dt 
= \frac{1}{2} \int_0^{\infty} e^{-st} \, dt 
- \frac{1}{2} \int_0^{\infty} e^{-st} \cos(2\omega t) \, dt
\]

Das erste Integral lässt sich dann folgendermaßen auflösen:

\[
\int_0^{\infty} e^{-st} \, dt = \left[ \frac{e^{-st}}{-s} \right]_0^{\infty} = \frac{1}{s}, \quad \text{für } \operatorname{Re}(s) > 0
\]

Um das zweite Integral aufzulösen nutzen wir die eulersche Formel, welche unseren Cosinusterm über eulersche Funktionen darstellbar macht:

\[
\cos(2\omega t) = \frac{e^{i 2\omega t} + e^{-i 2\omega t}}{2}
\]

Setzen wir diesen Term nun wieder in unser aufgestelltes zweites Integral ein, könne wir dieses wieder aufteilen:

\[
\begin{aligned}
\int_0^{\infty} e^{-st} \cos(2\omega t) \, dt 
&= \int_0^{\infty} e^{-st} \cdot \frac{e^{i 2\omega t} + e^{-i 2\omega t}}{2} \, dt \\
&= \frac{1}{2} \int_0^{\infty} \left( e^{-st + i 2\omega t} + e^{-st - i 2\omega t} \right) dt \\
&= \frac{1}{2} \left( \int_0^{\infty} e^{-t(s - i 2\omega)} dt + \int_0^{\infty} e^{-t(s + i 2\omega)} dt \right)
\end{aligned}
\]

Die Stammfunktionen der beiden Funktionen lassen sich darstellen zu:

\[
\int_0^{\infty} e^{-t(s \pm i 2\omega)} \, dt 
= \left[ \frac{-1}{s \pm i 2\omega} \, e^{-t(s \pm i 2\omega)} \right]_0^{\infty}
\]

Somit können wir das Integral auflösen:

\[
\int_0^{\infty} e^{-t(s \pm i 2\omega)} \, dt 
= \left[ \frac{-1}{s \pm i 2\omega} \, e^{-t(s \pm i 2\omega)} \right]_0^{\infty} 
= \frac{1}{s \pm i 2\omega}
\]

Somit ergibt sich für das zweite Integral:

\[
\int_0^{\infty} e^{-st} \cos(2\omega t) \, dt 
= \frac{1}{2} \left( \frac{1}{s - i 2\omega} + \frac{1}{s + i 2\omega} \right)
= \frac{1}{2} \cdot \frac{2s}{s^2 + 4\omega^2}
= \frac{s}{s^2 + 4\omega^2}
\]

Setzen wir das erste und zweite Integral nun wieder zusammen, bekommen wir den Term:

\[
\mathcal{L} \{ \sin^2(\omega t) \}
= \frac{1}{2} \cdot \left( \frac{1}{s} - \frac{s}{s^2 + 4\omega^2} \right)
\]

Die Terme müssen, um auf die Form in der Tabelle zu kommen, auf den gleichen Nenner gebracht werden:

Erster Term:

\[
\frac{1}{s} = \frac{s^2 + 4\omega^2}{s(s^2 + 4\omega^2)}
\]

Zweiter Term:

\[
\frac{s}{s^2 + 4\omega^2} 
= \frac{s \cdot s}{s(s^2 + 4\omega^2)} 
= \frac{s^2}{s(s^2 + 4\omega^2)}
\]

Subtrahiert man die Terme und multipliziert man mit \(\frac{1}{2}\), dann erhält man die gesuchte Form:

\[
\mathcal{L} \{ \sin^2(\omega t) \} = \frac{s}{2\omega^2 (s^2 + 4\omega^2)}
\]


